% 原 PPT 第 24 页；共 8 个显示状态。
\begin{frame}<1-8>[t]{\resizebox{94mm}{!}{题型一：求函数的解析式——代入法}}
\begin{campuscanvas}
\only<1|handout:0>{\node[anchor=north west,inner sep=0] at (25.600,116) {\includegraphics[page=103,width=153.6000mm]{source-content/content-trex.pdf}};}
\only<2|handout:0>{\node[anchor=north west,inner sep=0] at (25.600,116) {\includegraphics[page=104,width=153.6000mm]{source-content/content-trex.pdf}};}
\only<3|handout:0>{\node[anchor=north west,inner sep=0] at (25.600,116) {\includegraphics[page=105,width=153.6000mm]{source-content/content-trex.pdf}};}
\only<4|handout:0>{\node[anchor=north west,inner sep=0] at (25.600,116) {\includegraphics[page=106,width=153.6000mm]{source-content/content-trex.pdf}};}
\only<5|handout:0>{\node[anchor=north west,inner sep=0] at (25.600,116) {\includegraphics[page=107,width=153.6000mm]{source-content/content-trex.pdf}};}
\only<6|handout:0>{\node[anchor=north west,inner sep=0] at (25.600,116) {\includegraphics[page=108,width=153.6000mm]{source-content/content-trex.pdf}};}
\only<7|handout:0>{\node[anchor=north west,inner sep=0] at (25.600,116) {\includegraphics[page=109,width=153.6000mm]{source-content/content-trex.pdf}};}
\only<8|handout:1>{\node[anchor=north west,inner sep=0] at (25.600,116) {\includegraphics[page=110,width=153.6000mm]{source-content/content-trex.pdf}};}
\end{campuscanvas}
\end{frame}
